\section{Feature Definitions}\label{app:features}
\begin{table}[!htbp]
\caption{The 31 fixed digital features. Counts refer to scalar coordinates.
Texture measurements do not measure physical paint relief or authenticate brushwork.
Exact scales and definitions appear in Appendix~\ref{app:features}.}
\label{tab:features}

\centering\small
\begin{tabularx}{\linewidth}{@{}>{\raggedright\arraybackslash}p{.15\linewidth}>{\raggedright\arraybackslash}p{.35\linewidth}>{\raggedright\arraybackslash}X@{}}
\toprule
Family & Measurements & Interpretation\\
\midrule
Color (11) & Lightness and chroma median/IQR; chromatic fraction; hue
concentration and entropy; color differences at three lags and their slope &
Brightness, color strength and diversity, and how rapidly colors change over space.\\[4pt]
Spatial (8) & Spectral slope, residual and anisotropy; orientation entropy;
horizontal--vertical balance; gradient median/IQR; quadrant divergence &
Coarse versus fine organization, dominant directions, edge strength and regional
variation in orientation.\\[4pt]
Texture (12) & Four wavelet energies, slope and curvature; three local-binary-pattern
entropies; three local coefficients of variation &
Detail across scales, diversity of local patterns and local luminance contrast.\\
\bottomrule
\end{tabularx}


\end{table}

Color coordinates use CIELAB under D65 and the $2^\circ$ observer. Lightness
and chroma each contribute a median and IQR. Chromatic fraction is the share
of pixels with $C^*\geq5$. Hue concentration is the magnitude of the
chroma-weighted circular mean; hue entropy uses 24 equally spaced bins and is
normalized by their log count. Both use chroma weights and are zero when fewer
than 1\% of pixels are chromatic. At each of three lags (1\%, 4\% and 16\% of
the short side), the median is taken over pooled horizontal and vertical
CIEDE2000 differences. A log--log slope summarizes their
scale dependence \citep{sharma2005ciede}.

Spatial and texture coordinates use linear-sRGB luminance. A Tukey window
(parameter .1) precedes Fourier analysis. Thirty-two logarithmic radial bins
cover 4--128 cycles per short side; a Theil--Sen line gives spectral slope and
its root-mean-square residual. Power-weighted axial concentration measures
anisotropy. Scharr derivatives supply gradient median and IQR, an 18-bin
magnitude-weighted orientation entropy, and horizontal--vertical balance. Mean
Jensen--Shannon divergence of quadrant orientation histograms from the global
histogram measures regional variation.

Four log mean-squared detail energies come from a four-level undecimated
\texttt{db2} wavelet transform, followed by their linear slope and mean second
difference \citep{mallat1989wavelet}. Uniform local-binary-pattern entropy uses
$(P,R)=(8,1),(16,2),(32,4)$ on quantized luminance
\citep{ojala2002lbp}. Three median local coefficients of variation use window
widths obtained by rounding 1\%, 4\% and 16\% of the short side and increasing
even widths to the next odd integer, with a $10^{-6}$ mean floor.
The fixed implementation records boundary handling and stabilizing
constants. The separate scaler contains 101 Monet, 36 Sisley, 48 Pissarro and
36 C\'ezanne works. Equal painter mass is used to calculate each weighted
median and IQR.

\section{Estimator Details and Sensitivity Analyses}\label{app:estimation}
Let $\theta_{msa}=\mathbb E[d_{msak}]$ and
$d_{msak}=\theta_{msa}+\eta_{msak}$. If repeat errors have zero mean and
$\mathbb E\langle\eta_{msa1},\eta_{msa2}\rangle=0$, then
\[
 \mathbb E[\widehat D_{ms}]
 =H^{-1}\sum_a\|\theta_{msa}-r_a\|^2.
\]
Centering induces dependence between artists within a repeat, but independent
repeat blocks preserve the vanishing cross-repeat term. A common response has
$\theta_{msa}=0$ and hence expected $D=1$, even if it shifts the entire named
collection toward the references. A correctly aligned but rescaled response
$\theta_{msa}=b r_a$ gives expected $D=(b-1)^2$. Orthogonal deviations add their
squared magnitude. This identity explains why stronger painter contrasts need not have lower
reference-relative squared error.

The aggregate error $D_{\rm aggregate}$ applies the cross-product estimator to scene-averaged
contrasts; Equation~\ref{eq:scene-decomposition} relates it to the primary,
equally scene-weighted error. Painters also receive equal weight. Their
additive contributions sum to $D$, but centering couples the four conditions,
so these contributions are not standalone artist-fidelity scores.

For a model contrast, compute one error difference per complete paired scene.
Its standard error is the sample standard deviation of those differences divided
by $\sqrt n$. The simultaneous half-width is
$t_{n-1,1-.05/(2\cdot21)}\,\mathrm{SE}$. Six aligned-response intervals use the same
family adjustment. Additional nominal intervals and 5,000 paired-scene bootstrap
draws are descriptive checks, not alternative routes to significance. Reference
sensitivity draws 1,000 bootstrap samples independently within each painter,
recomputes reference means and $H$, and holds generated vectors and scaling
fixed. It describes reference-panel dependence rather than all uncertainty
jointly.

Declared before inspecting the new feature outcomes, the four-class reference
sensitivity uses a fixed lexicon applied to source titles: water-organized, built-place-organized, route-organized, and open or
wooded land. Each class mean receives one-quarter weight within a painter.
Resampling within artist/class strata conditions on those labels and does not
account for classification error. The older controlled naming panel used a
separate three-class visual coding scheme, described below.

For the shared-effect diagnostic, average each arm over scenes within repeat.
With named-minus-free differences $h_{ak}$ and $c_k=\frac14\sum_a h_{ak}$,
common and specific squared magnitudes are $4\langle c_1,c_2\rangle$ and
$\sum_a\langle h_{a1}-c_1,h_{a2}-c_2\rangle$. Their common fraction is reported
when the total is positive; corrected components may be negative. The within-scene
version computes both squared-change components separately in each scene,
averages each over scenes, then takes their ratio. Its shares range from 88.6\%
(GPT Image 1) to 97.2\% (FLUX).
Let $g_k$ be the generic-minus-free shift. The ordinary repeat-averaged cosine
shares a noisy free estimate in both vectors, adding its error covariance trace
to the expected numerator. The corrected cosine reported in
Section~\ref{sec:shared-change} instead uses cross-repeat products,
$[\langle g_1,c_2\rangle+\langle g_2,c_1\rangle]/
[2\sqrt{\langle g_1,g_2\rangle\langle c_1,c_2\rangle}]$ when both norm squares
are positive. The ratio itself is not unbiased or constrained to $[-1,1]$.
Primary $D$ uses only named arms and is unaffected by shared-free-control noise.

\subsection{Additional Score Identities}\label{app:score-details}
Under a shared affine transformation $z\mapsto Az+b$, centering removes
the additive term $b$ but leaves the transformed contrast $Ad$.

The estimator has an exact decomposition. Write
$e_{msak}=d_{msak}-\widehat\beta_{msk}r_a$. Then
\begin{equation}
 \widehat D_{ms}=
 (\widehat\beta_{ms1}-1)(\widehat\beta_{ms2}-1)
 +\frac1H\sum_a\langle e_{msa1},e_{msa2}\rangle.
 \label{eq:decomposition}
\end{equation}
The first term measures reference-aligned amplitude error. The residual is
orthogonal to the single stacked reference configuration; it is not necessarily
outside the feature-space span of the reference means, much less outside all
authentic artistic variation. Both terms can be negative in finite samples.

Let $\theta_s$ stack the four conditional mean contrasts and
$\bar\theta=S^{-1}\sum_s\theta_s$. The population target decomposes as
\begin{equation}
 D_{\rm scene}=\frac{\|\bar\theta-r\|^2}{H}
 +\frac1{SH}\sum_s\|\theta_s-\bar\theta\|^2
 =D_{\rm aggregate}+V_{\rm scene}.
 \label{eq:scene-decomposition}
\end{equation}
To separate contrast strength from direction, define
\begin{equation}
 \widehat Q_m=\frac1{SH}\sum_{s,a}\langle d_{msa1},d_{msa2}\rangle,
 \qquad \widehat D_m=1-2\widehat\beta_m+\widehat Q_m.
 \label{eq:magnitude}
\end{equation}
For positive $\widehat Q_m$, $\widehat\beta_m/\sqrt{\widehat Q_m}$ is a
noise-corrected alignment summary; finite estimates need not lie in $[-1,1]$.
Scaling every centered generated contrast by $c$ gives
$\widehat D_m(c)=1-2c\widehat\beta_m+c^2\widehat Q_m$.
We fit $c=\max(0,\widehat\beta/\widehat Q)$ on the other 13 scenes and evaluate
it on the held-out scene. The fitting criterion is available only when its
training $\widehat Q$ is positive, as it is in every observed fold. This
counterfactual concerns feature vectors, not an implemented image transformation.

\subsection{Paired Comparisons and Declared Sensitivities}\label{app:declared-sensitivities}
We fit the scene fixed-effects regression
\begin{equation}
 \widehat D_{ms}=\alpha_s+\gamma_m+\varepsilon_{ms},
 \qquad\gamma_{\text{GPT Image 2}}=0.
 \label{eq:regression}
\end{equation}
In the balanced panel, coefficient differences equal the paired scene means
used for inference. Changing the regression baseline leaves these differences
unchanged; Table~\ref{tab:specificity-models} presents them relative to FLUX.
Declared sensitivities delete each scene, resample scenes
and reference works, examine feature families and omit texture. Central-square
windows control aspect ratio by discarding peripheral content in generated and
reference images. FLUX retains the lowest uncalibrated error after every scene
deletion ($D=.754$--$.852$) and with square windows ($D=.765$; next-lowest
Nano Banana 2, $1.024$).

\subsection{Distributional Proximity and Repeat Spread}
For painter $a$, energy discrepancy compares its generated collection $Y_a$
with reference collection $X_a$:
\begin{equation}
 \E(X,Y)=\frac{2}{n_Xn_Y}\sum_{i,j}\|x_i-y_j\|
 -\frac1{n_X^2}\sum_{i,i'}\|x_i-x_{i'}\|
 -\frac1{n_Y^2}\sum_{j,j'}\|y_j-y_{j'}\|.
 \label{eq:energy}
\end{equation}
Its finite-sample expectation depends on within-collection spread; equal
generated sample sizes do not eliminate comparison bias. Appendix~\ref{app:diagnostics}
gives the correction for the repeated-scene design. Within-scene repeat spread,
$\frac{1}{2S}\sum_s\|z_{msa1}-z_{msa2}\|^2$, is the mean sample-variance
trace of the two outputs. It separates repeat dispersion from variation across
briefs.


\section{Post-result Target Diagnostics}\label{app:diagnostics}
These retrospective analyses retain the completed 1,008 generated images and
the original inferential family. The initial diagnostics use unchanged reference
vectors and scaling; the subsequent source audit separately varies painting
regions, development scaling and content labels.

\subsection{Content-specific Targets and Genuine-painting Controls}
The recorded reference classes are water-organized, built-place-organized,
route-organized, and open or wooded land, assigned by a fixed source-title
lexicon. Counts in that order are $(191,67,8,31)$ for Monet, $(52,29,9,16)$ for
Sisley, $(28,77,12,24)$ for Pissarro and $(31,28,12,34)$ for C\'ezanne.
Class-specific means are centered across painters just as in the pooled target.
Their mean squared departure from the pooled target is $.438H$, a description
of observed reference means that includes finite-sample error.

For the generated comparison, the four water, four built and three land briefs
use the corresponding class means; the three mixed briefs are excluded.
We average the class-specific cross-product error over these $S=11$ scenes and divide by
$H_c=S^{-1}\sum_s\sum_a\|r_{a,c(s)}\|^2=7.413$, compared with pooled $H=5.915$.
Each column of Table~\ref{tab:review-targets} therefore has its own denominator.
The unchanged 11-scene pooled comparison isolates the scene-subset choice.
These title-derived strata have 16--191 reference works per artist/class among
the three included classes, but do not supply visually matched historical scenes.

For each of 1,000 genuine-painting controls, works in each artist/class cell
are randomly divided into disjoint halves, with the smaller half used for
target estimation. A pooled control samples two independent works with replacement
from each painter's held-out half for each of 14 pseudo-scenes. A class-stratified
control samples within the held-out class for each of the 11 water, built or
land pseudo-scenes. It is compared with both pooled and class-specific means
from the training halves, using each target's own squared magnitude.
The random seed is 2026091201. Sampling can repeat a held-out work, but no
held-out work is used to estimate that split's target. The development scaler
is fixed throughout.

To separate sparse sampling from changes in the comparison pools, we also
calculate each split's expected score from its exact held-out means, keeping
the training target and normalizer fixed (Table~\ref{tab:painting-control-ranges}).

\begin{table}[htbp]
\caption{Central 95\% ranges of genuine-painting control scores across 1,000
splits. Exact means integrate out sampling within each split; these finite-pool
ranges are not confidence intervals.}
\label{tab:painting-control-ranges}

\centering\small
\begin{tabular}{llcc}
\toprule
Control sampling & Reference target & Two draws & Exact held-out means\\
\midrule
Pooled & Pooled & $[-.758,1.380]$ & $[.135,.428]$\\
Class-stratified & Pooled & $[-.554,2.400]$ & $[.544,.993]$\\
Class-stratified & Class-specific & $[-.248,1.946]$ & $[.517,.954]$\\
\bottomrule
\end{tabular}


\end{table}

For pooled controls, SD falls from $.561$ to $.076$ after integrating out
sampling. These controls still condition on imperfect title classes and a
single source capture per work; they do not calibrate artistic fidelity.

\begin{table}[htbp]
\caption{Descriptive error under alternative targets and metrics. The class target uses title-derived water, built and land strata; the three mixed briefs are excluded. Each column uses its own reference normalizer; absolute values across columns are not directly comparable. Covariance is fitted only on development works, with fixed 50\% shrinkage.}
\label{tab:review-targets}
\centering\small
\begin{tabular}{@{}lrrrr@{}}\toprule
 & \multicolumn{2}{c}{11 common briefs} & \multicolumn{2}{c}{14 briefs}\\
Model & Pooled target & Class target & Equal family & Covariance\\\midrule
GPT Image 1 & 1.557 & 1.366 & 1.551 & 2.103\\
GPT Image 2 & 1.099 & 1.130 & 1.243 & 1.511\\
2.5 Flare & 1.726 & 1.539 & 1.765 & 1.636\\
2.5 Sunburst & 1.633 & 1.525 & 1.680 & 1.617\\
Nano Banana 2 & 0.987 & 0.948 & 1.203 & 1.385\\
FLUX.2 Max & 0.750 & 0.786 & 0.808 & 0.928\\
\bottomrule\end{tabular}

\end{table}

\FloatBarrier

\subsection{Magnitude, Artist Coverage and Coordinate Weighting}\label{app:magnitude-coverage}
The cross-repeat version of Equation~\ref{eq:scene-decomposition} is exact for
the retained vectors. Its variation term replaces squared deviations with
cross-products between scene-centered repeats. The term can be negative in
finite samples, although its independent-repeat population target is nonnegative.
Scalar calibration fits only the other 13 scenes. All fold denominators $Q$
are positive. GPT Image 2's fitted scalars range from .438 to .459 and FLUX's
from .598 to .694; all fold-specific scores are retained. No calibration
parameter is selected by minimizing its own held-out score. The subsequent
deletion check refits the entire procedure on
each remaining 13-scene panel, rather than treating overlapping folds as
independent replications. Its mean GPT Image 2-minus-FLUX error differences
range from $-.189$ to $-.137$.

\begin{table}[htbp]
\caption{Descriptive magnitude and alignment of painter contrasts. $Q=1$ equals the reference squared size, and $D=1-2\beta+Q$. The alignment ratio can fall outside $[-1,1]$ in finite samples.}
\label{tab:review-alignment}
\centering\small
\begin{tabular}{@{}lrr@{}}\toprule
Model & Squared contrast magnitude $Q$ & Alignment ratio $\beta/\sqrt Q$\\
\midrule
GPT Image 1 & 2.449 & 0.600\\
GPT Image 2 & 2.223 & 0.670\\
2.5 Flare & 2.281 & 0.511\\
2.5 Sunburst & 2.289 & 0.545\\
Nano Banana 2 & 0.994 & 0.444\\
FLUX.2 Max & 0.741 & 0.546\\
\bottomrule\end{tabular}

\end{table}


For artist pairs, both generated and reference contrasts are direct differences
between the two artists; the denominator is the squared reference-pair distance.
Leaving one artist out recenters the remaining three generated and reference
conditions. Reference singular components are the three rank-one terms of the
centered four-by-31 reference matrix: each combines an artist contrast with a
feature direction, rather than representing one artist. Component aligned
responses project onto each term separately.

\begin{table}[htbp]
\caption{Descriptive artist coverage of the aligned response. Reference components account for 66.3\%, 20.6\% and 13.0\% of centroid variation. Each omission recomputes the target and recenters the other three artists.}
\label{tab:review-components}
\centering\small
\begin{tabular}{@{}lrrr|rrrr@{}}\toprule
 & \multicolumn{3}{c}{Component aligned response} & \multicolumn{4}{c}{Aligned response after omitting}\\
Model & First & Second & Third & Monet & Sisley & Pissarro & C\'ezanne\\\midrule
GPT Image 1 & 1.248 & 0.310 & 0.356 & 1.140 & 1.072 & 0.840 & 0.389\\
GPT Image 2 & 1.156 & 0.617 & 0.800 & 1.195 & 0.972 & 0.992 & 0.651\\
2.5 Flare & 1.022 & 0.211 & 0.383 & 1.034 & 0.734 & 0.802 & 0.222\\
2.5 Sunburst & 1.082 & 0.250 & 0.417 & 1.116 & 0.840 & 0.762 & 0.284\\
Nano Banana 2 & 0.513 & 0.371 & 0.195 & 0.569 & 0.446 & 0.388 & 0.276\\
FLUX.2 Max & 0.666 & 0.021 & 0.183 & 0.539 & 0.511 & 0.527 & 0.096\\
\bottomrule\end{tabular}

\end{table}


Table~\ref{tab:review-components} shows that aggregate positive alignment is
uneven across those components.

Equal-family weighting assigns each coordinate inverse family size (11, 8 or 12).
The covariance sensitivity uses 221 development works with equal painter mass.
For their weighted covariance $\Sigma$, the metric is
$[.5\Sigma+.5\operatorname{tr}(\Sigma)I/31]^{-1}$, with fixed shrinkage and no
tuning on evaluated images. These weighting choices are not independently
validated artistic metrics.

\begin{samepage}
FLUX retains the lowest point estimate under all
31 single-coordinate deletions; the report retains every coordinate contribution,
including negative terms.
Without texture, Flare and Sunburst still have error estimates above the
no-contrast value of one (1.580 and 1.604).
\end{samepage}

\subsection{A Design-aware Energy Correction}
The primary energy values use empirical distributions, including their zero
diagonal distances. To estimate energy between a fixed empirical reference
and the equally weighted mixture of scene-conditional output distributions,
cross-scene generated pairs already have the required independent draws.
Only the same-scene generated term needs adjustment. With two repeats and
$S$ scenes, the descriptive correction is
\[
 \widehat{\mathcal E}_{\rm strata}
 =\widehat{\mathcal E}_{\rm empirical}
 -\frac{1}{2S^2}\sum_s\|y_{s1}-y_{s2}\|.
\]
The reference within-distance remains empirical because this target conditions
on the finite reference panel. A pooled IID $n/(n-1)$ multiplier would also
alter cross-scene terms and estimate a different target. This correction retains
the 21 negative named-minus-generic differences. It does not solve unequal
historical and generated content mixtures, and no new energy tests are performed.

\subsection{Uncertainty and Synthetic Coverage Checks}
The Student procedure uses 14 scene summaries. For independent scene errors
with fixed conditional means $\delta_s$, its estimated variance has expectation
\[
 \mathbb E[s_\delta^2/S]
 =\operatorname{Var}(\bar{\widehat\delta})
 +\frac{1}{S(S-1)}\sum_s(\delta_s-\bar\delta)^2.
\]
It therefore includes real between-scene heterogeneity as well as request noise.
This explains why it is not a direct estimate of fresh-request variance for
exactly the fixed panel. It also does not provide an exact coverage guarantee:
non-Gaussian cross-products and dependence can affect the Student approximation.

The initial 5,000-run scenarios (seed 2026091202) use the six configurations,
14 scenes, two repeats and 21-endpoint family, with unit reference magnitude.
Gaussian, variance-standardized $t_3$/heteroskedastic and fixed scene-interaction
errors give simultaneous coverage 95.96\%, 97.42\% and 99.68\%. Their independent
centered noise power is $.5$. Allocating half the noise to a model-specific
state shared by all scenes and both repeats reduces coverage to 0.04\%, with
mean corrected-error bias near $.25$.

A subsequent sweep uses independent Gaussian errors at powers $.5$, $1.5$ and
$3$, with either isotropic or rank-three covariance aligned with the reference
components (5,000 runs per case). Coverage ranges from 95.06\% to 96.80\%, with
Monte Carlo SE $.25$--$.31$ percentage points. For comparison, directly centered
repeat differences estimate observed powers of $.424$--$2.596$ relative to $H$.
These stress models probe coverage approximations, not actual service independence;
no significance threshold is selected from them.

\subsection{Source-quality Sensitivity}\label{app:reference-quality}
AI assistants inspect all 649 reference and 221 development reproductions
using contact sheets and enlarged peripheral regions, recording removable
calibration targets, frames, margins and labels. Uncertain boundaries remain
unchanged. Reference content is visually classified before comparison with the
title label; mixed or unclear cases retain that label.

The audit selects crops for 90 reference and 41 development images, including
ten calibration-strip cases; 43 uncertain boundaries remain unchanged. Visual
classes differ from 230 title assignments, while 42 mixed or unclear cases
retain their original labels. These AI-coded sensitivities did not undergo human verification and are not
independent ground truth.

Audited rectangular crops follow the original orientation and ICC rules, before
aspect-preserving Lanczos resizing to a 512-pixel short side without upsampling.
All 31 features and all works are retained. We compare corrected references with
original scaling, then refit development IQR scaling and consistently transform
generated vectors. Label sensitivity uses the same 11 briefs. Versioned records
preserve audit decisions, new vectors and original results; model scores do not
select corrections.
All six aligned-response intervals remain above zero with either scaler.
With original scaling, corrected regions give FLUX $D=0.832$.
After refitting, (uncalibrated, held-out calibrated) errors are:
GPT Image 1 $(1.631,0.693)$, GPT Image 2 $(1.259,0.579)$,
Flare $(1.737,0.760)$, Sunburst $(1.577,0.711)$, Nano Banana 2 $(1.121,0.846)$
and FLUX $(0.872,0.762)$. FLUX and GPT Image 2 retain the respective minima.

\section{Image Selection and Provenance}\label{app:images}

\begin{figure}[p]
\centering\includegraphics[width=\linewidth]{specificity_original_generated.pdf}
\caption{One scene, six prompt conditions and six generators. The requested scene
is a broad river bending past a gravel bank, three trees on the opposite bank,
and a low hill under an overcast sky. In each model row, compare the four named
outputs with the artist-free and generic oil-painting controls on the right,
then compare the painter names with one another. The top row shows audited
reference regions for display; primary scores use full frames. These works were
neither generator inputs nor matched compositions. All generated examples use the first brief and first
repeat; Appendix~\ref{app:images} gives the full prompt, selection rule and provenance.}
\label{fig:original-generated}
\end{figure}

The generated examples use the first declared brief and first repeat, selected
after analysis without appearance- or score-based curation. The brief reads:
``A broad river bends past a gravel bank. Three tall trees stand on the far bank,
with a low hill behind them under a pale overcast sky.'' Every prompt ends with
``No text or frame.'' Only the clauses in Section~\ref{sec:design} change.

References in Figure~\ref{fig:original-generated} are the lexicographically first
measured work IDs per painter classified as water-organized in the completed
audit, with recorded public-domain status. The Sisley and Pissarro reproductions
use audited crops removing calibration strips; Monet and C\'ezanne retain their
full regions. The manifest records image identities, provenance, crop boxes
and exact prompts. The paintings are not matched to the generated brief.


\FloatBarrier


\section{Complete Model and Painter Summaries}\label{app:results}
Tables~\ref{tab:all-specificity-pairs}--\ref{tab:artist-distributions} and
Figures~\ref{fig:specificity-diagnostics}--\ref{fig:specificity-distributions}
summarize model comparisons, error components and collection-level proximity
and spread.
\emph{Energy discrepancy} combines between-collection distances and internal
spread; lower values mean greater proximity (Appendices~\ref{app:estimation}
and~\ref{app:diagnostics}). Naming lowers energy relative to the generic clause in 21 of 24
model--painter cells, including the correction for repeated scenes.
All 24 named collections have smaller total feature variance than their historical
counterparts. These descriptive comparisons use 28 named or generic outputs
against the recorded reference panel, with different subject mixtures; they do
not establish agreement of painter contrasts.
\begin{table}[htbp]
\caption{All 15 paired model comparisons in the prespecified inferential family. Intervals are approximate 95\% simultaneous intervals adjusted across all 21 endpoints. A negative contrast favors model A on the stated feature-geometry target.}
\label{tab:all-specificity-pairs}
\centering\small
\begin{tabular}{llr}\toprule
Model A & Model B & $D_A-D_B$ [simultaneous interval]\\\midrule
GPT Image 1 & GPT Image 2 & $0.346\;[-0.400,1.093]$\\
GPT Image 1 & 2.5 Flare & $-0.165\;[-1.027,0.697]$\\
GPT Image 1 & 2.5 Sunburst & $-0.069\;[-1.181,1.043]$\\
GPT Image 1 & Nano Banana 2 & $0.463\;[-0.602,1.528]$\\
GPT Image 1 & FLUX.2 Max & $0.771\;[-0.037,1.580]$\\
GPT Image 2 & 2.5 Flare & $-0.512\;[-1.150,0.126]$\\
GPT Image 2 & 2.5 Sunburst & $-0.415\;[-1.181,0.350]$\\
GPT Image 2 & Nano Banana 2 & $0.116\;[-0.922,1.155]$\\
GPT Image 2 & FLUX.2 Max & $0.425\;[-0.161,1.011]$\\
2.5 Flare & 2.5 Sunburst & $0.096\;[-0.286,0.478]$\\
2.5 Flare & Nano Banana 2 & $0.628\;[-0.283,1.540]$\\
2.5 Flare & FLUX.2 Max & $0.937\;[0.256,1.618]$\\
2.5 Sunburst & Nano Banana 2 & $0.532\;[-0.436,1.500]$\\
2.5 Sunburst & FLUX.2 Max & $0.840\;[0.058,1.623]$\\
Nano Banana 2 & FLUX.2 Max & $0.309\;[-0.847,1.464]$\\
\bottomrule\end{tabular}

\end{table}

\begin{table}[htbp]
\caption{Complementary empirical distribution summaries. Energy change is named-minus-generic reference energy; variance ratios are generated-to-reference. Negative changes favor naming; ratios below one indicate narrower generated feature distributions.}
\label{tab:artist-distributions}
\centering\small
\begin{tabular}{llrrrr}\toprule
Model & Measure & Monet & Sisley & Pissarro & C\'ezanne\\\midrule
GPT Image 1 & Energy change & $-2.515$ & $-2.739$ & $-2.628$ & $-1.336$\\
& Variance ratio & $0.479$ & $0.694$ & $0.603$ & $0.598$\\
\addlinespace[2pt]
GPT Image 2 & Energy change & $-0.602$ & $-1.092$ & $-0.195$ & $-0.373$\\
& Variance ratio & $0.283$ & $0.312$ & $0.303$ & $0.303$\\
\addlinespace[2pt]
2.5 Flare & Energy change & $0.962$ & $-0.008$ & $0.429$ & $-1.013$\\
& Variance ratio & $0.269$ & $0.371$ & $0.321$ & $0.275$\\
\addlinespace[2pt]
2.5 Sunburst & Energy change & $0.267$ & $-0.718$ & $-0.091$ & $-1.653$\\
& Variance ratio & $0.248$ & $0.274$ & $0.273$ & $0.318$\\
\addlinespace[2pt]
Nano Banana 2 & Energy change & $-0.513$ & $-1.381$ & $-0.741$ & $-1.102$\\
& Variance ratio & $0.829$ & $0.708$ & $0.572$ & $0.877$\\
\addlinespace[2pt]
FLUX.2 Max & Energy change & $-0.707$ & $-0.886$ & $-1.080$ & $-1.431$\\
& Variance ratio & $0.340$ & $0.416$ & $0.570$ & $0.622$\\
\bottomrule\end{tabular}

\end{table}

\FloatBarrier

\begin{figure}[tbp]
\centering\includegraphics[width=\linewidth]{specificity_diagnostics.pdf}
\caption{Descriptive summaries for the original reference target.
(a) Amplitude error and the orthogonal residual---differences perpendicular to
the overall labeled reference pattern---sum to $D$ (Equation~\ref{eq:decomposition}).
(b) Generated/reference variance ratios for every model and painter. Total
spread includes variation across fixed briefs and is not a mode-collapse measure.}
\label{fig:specificity-diagnostics}
\end{figure}
\begin{figure}[tbp]
\centering\includegraphics[width=.92\linewidth]{specificity_sunburst_distributions.pdf}
\caption{GPT Image 2.5 Sunburst's named outputs appear more concentrated than
the reference works in all four projections. This model illustration was selected before the
primary feature outcomes were inspected. Each panel uses principal components
fitted only to that painter's reference works, so axes differ between painters.
The 14 generated briefs and two repeats have a different content distribution
from the historical panel. Projections are descriptive; numerical summaries
use all 31 coordinates.}
\label{fig:specificity-distributions}
\end{figure}
\FloatBarrier

\section{Supporting Interventions}\label{sec:why}\label{app:cohorts}
These earlier cohorts address related but different questions
(Table~\ref{tab:cohorts}). They are not pooled with the six-model experiment
or treated as independent-investigator replications. Historical model labels
identify recorded request configurations, not verified underlying models
(Appendix~\ref{app:provenance}).

\begin{table}[htbp]
\caption{Supporting evidence is kept separate from the new six-model experiment.
Different prompts, panels and finite-sample targets do not form one pooled test.}
\label{tab:cohorts}

\centering\small
\begin{tabularx}{\linewidth}{@{}>{\raggedright\arraybackslash}p{.22\linewidth}>{\raggedright\arraybackslash}p{.27\linewidth}>{\raggedright\arraybackslash}X@{}}
\toprule
Cohort & Comparison & Role in this paper\\
\midrule
Four-painter exploration & 649 references; 1,920 generated outputs &
Describes reference/generated separation and motivates the controlled test.\\[4pt]
Controlled naming & 1,006 outputs; 38 Monet and 32 C\'ezanne references &
Six named/free and two detailed/short prompt contrasts; subsequent transformation
and conditional-variance diagnostics.\\[4pt]
Palette interventions & Two separate 192-output collections &
Named-minus-generic chroma response; no pooling across collections.\\[4pt]
Generic-clause comparison & 96 C\'ezanne/generic outputs on 24 new scenes &
Tests a proximity gain beyond a generic painting clause.\\[4pt]
Fixed-map transfer & 240 FLUX/C\'ezanne outputs on 12 new scenes &
Prospective counterexample to the historical opposing energy/prediction ordering.\\
\bottomrule
\end{tabularx}


\end{table}

\subsection{Proximity and Spread under Naming}
The four-painter exploration used 1,536 named and 384 artist-free outputs,
three prompt methods and 16 scenes. Its 24 model--painter--method cells had
total variance ratios .206--.376 and grouped generated/reference discrimination accuracy
.940--.983. Pooled named energy medians for Monet, Sisley, Pissarro and C\'ezanne
were 2.068, 1.678, 1.975 and 2.142; corresponding artist-free medians were
1.732, 1.878, 1.854 and 1.657. Only Sisley had a lower pooled named median.
Matched-size real/real medians were .193, .215, .164 and .199. These post-result
summaries describe that exploration, not a general naming effect.

The later controlled panel selected Monet and C\'ezanne after exploration.
It crossed 24 scenes with three models, two clauses and three repeats; 142
short-scene outputs supplied two prompt-length comparisons. Reference panels
contained 38 Monet and 32 C\'ezanne works, with water/built/land masses based on
maintainer-run LLM visual coding. Named-minus-free energy changes were
$-.846,-1.818$ for Nano Banana 2, $-.625,-.896$ for FLUX, and $-.234,-.056$ for
GPT Image 2 (Monet, C\'ezanne). The first four null hypotheses were rejected in
the original family of eight. A separate 72-output FLUX collection repeated both directions
($-.695,-.952$) with rejection in its own four-endpoint family.

A standalone 96-output GPT Image 2/C\'ezanne comparison on 24 new scenes reduced energy
relative to the generic clause by 1.195 ($p=.00001$, fixed $\alpha=.025$), with
named/generic total variance ratio .552. It randomized clauses within 48 two-position
blocks, used reference class masses $(3,11,18)/32$, and tested a sharp no-effect
null with 99,999 within-block sign draws and plus-one correction. This tests
clause effects, not equality of energy alone. A generic clause therefore did
not reproduce the full proximity gain in that panel, but proximity does not
establish correct four-artist contrasts. Likewise, lower total spread does
not establish loss of scene information: controlled FLUX/Monet held-repetition
scene retrieval rose from 44.4\% to 59.7\% under naming. This measures relative
feature distinguishability, not verified prompt adherence.

\subsection{Palette Responsiveness and a Contrary Transfer Result}
Both 192-output palette collections requested the historical GPT Image 2 alias,
using six scenes, four repeats, muted/vivid instructions and free, generic,
Monet and C\'ezanne clauses. Primary name-minus-generic chroma-response
interactions were unresolved (Monet, C\'ezanne):
$-.284,-.018$ in the first collection and $-.159,.037$ in the later one.
Their adjusted intervals spanned zero in the original two-interaction family and the later four-endpoint
family, which also included two FLUX naming contrasts.
Secondary generic-minus-free responses were $-.853$ and $-.891$, with nominal
95\% intervals excluding zero. Thus generic painting can attenuate measured
palette response without an identified additional name effect. Two palette
levels do not distinguish reduced gain from saturation or an operating-range
shift, and unresolved interactions do not demonstrate equivalence.

Earlier maps translated and optionally isotropically rescaled artist-free feature
vectors to predict named-prompt scene means, with fits based only on generated
images and evaluation on held-out scenes. Actual naming had lower reference energy than
translation plus isotropic scaling in all six model--painter averages. Adding
scale improved named-mean prediction in all six while worsening reference
energy in five. That ordering did not transfer to a prospective 240-output
FLUX/C\'ezanne comparison with 12 new scenes and ten repeats, using the original
fits unchanged. Translation/scaling improved both energy and repeat-corrected
prediction error relative to translation: $-.390\;[-.509,-.271]$ and
$-5.047\;[-7.880,-2.214]$, with approximate simultaneous jackknife intervals.
The achieved prediction-error half-width exceeded its planning criterion.
This contrary result is retained without extra sampling; it rules out treating
the historical opposing ordering as a universal mechanism. These named/free maps
differ from the retrospective calibration of painter contrasts in
Table~\ref{tab:review-calibration}.

\section{Collection Provenance and Computational Scope}\label{app:provenance}
Historical GPT Image 1 and GPT Image 2 labels identify the requested model
names. A later technical negative control showed that the earlier image
interface also accepted an invented image-model identifier. This does not show
that valid names select the same underlying model, but successful request
acceptance alone cannot authenticate their separation. Historical contrasts
are therefore interpreted as recorded request-group comparisons. The current
experiment uses explicit documented model identifiers and fixed routes;
it does not reuse historical images as verified examples of the new variants.
All current outputs decode to the requested square dimensions, but their
response bodies do not echo model or quality fields. Labels therefore denote
the documented request configurations, without independent checkpoint
attestation. Underlying checkpoints and training corpora are not available
for inspection.

The present experiment excludes technical probes and an earlier attempt closed
before feature measurement because model selection could not be qualified.
No image is shared between these collections. The earlier incomplete
generic-clause cohort is likewise not pooled with its complete 96-image
successor. Exact assignments, attempt outcomes and measurement exclusions are
retained in the versioned manifests; selection was not based on visual success.

\clearpage
\section{Projection of Labeled Painter Contrasts}\label{app:geometry}
In Figure~\ref{fig:specificity-geometry}, C\'ezanne remains visibly separated,
while generated Monet and Sisley means often lie close together. The M/S/P detail
boxes magnify the nearby positions. The axes combine the 31 standardized features, with directions fitted
only to the four reference means. Together they show 86.99\% of variation among
those means; Figure~\ref{fig:artist-pairs} checks all 31 features.

\begin{figure}[tbp]
\centering\includegraphics[width=\linewidth]{specificity_geometry.pdf}
\caption{Matching the labeled painter pattern. Crosses mark reference contrasts;
open circles average generated contrasts over all 14 scenes and both repeats.
Shorter matching-label lines indicate less projected mismatch;
faint dots show centered scene/repeat contrasts.
The outlined region is enlarged in the adjacent Monet (M), Sisley (S) and
Pissarro (P) detail at a common scale; C\'ezanne
remains in the overview. Each view uses common axes across models. Agreement
in these two projected axes does not establish agreement in all 31 features.}
\label{fig:specificity-geometry}
\end{figure}
